\documentclass[../main.tex]{subfiles}

\begin{document}

\section{Linear Algebra}

Linear algebra is a foundational building block of ML as majority
of the components are written in very high dimensional arrays
and tensors. Without exaggeration, training and using a model
can be considered as performing billions of linear algebra
operations. Therefore, it is encouraged to get familiar with
the notations as they will be used throughout the future notes.
That being said, because many of the contents were introduced
in detail in the linear algebra part of LibreNotebook, I will
keep it brief here for topics that were not introduced in the
previous part.

\subsection{Foundations}

Many mathematicians might not like the definition below, but I think
it is worth considering different perspectives of how we can see
vectors and matrices.

\begin{definition}{}{}
A 1-dimensional vector is simply called \textbf{vector}. A 2D vector
is called \textbf{matrix}. For vectors with 3+ dimensions, the term
\textbf{tensor} is used.
\end{definition}

Now scalars are not considered vectors because it is only defined
with magnitude. Geometrically speaking, the dimensions here can be
interpreted as the dimension that we usually say in science.
For example, a matrix can be plotted in a 2D space with $x$ and $y$
components. Similarly, a 3D tensor can be drawn in 3-dimensional
plane with $x$, $y$, and $z$ coordinates. Consider the following
example.

\begin{figure}[H]
\centering
\tdplotsetmaincoords{70}{110}
\begin{tikzpicture}[
        tdplot_main_coords,
        axis/.style={->,thick},
        scale=0.5
    ]
    \foreach \x in {0,1,...,7}
        \foreach \y in {0,1,...,7} {
			\draw[gray,very thin]
                (\x,0,0) -- (\x,7.5,0);
			\draw[gray,very thin] (0,\y,0) -- (7.5,\y,0);
		}
    \foreach \y in {0,1,...,7}
        \foreach \z in {0,1,...,7} {
			\draw[gray,very thin]
                (0,\y,0) -- (0,\y,7.5);
			\draw[gray,very thin] (0,0,\z) -- (0,7.5,\z);
		}
    \foreach \z in {0,1,...,7}
        \foreach \x in {0,1,...,7} {
			\draw[gray,very thin]
                (0,0,\z) -- (7.5,0,\z);
			\draw[gray,very thin] (\x,0,0) -- (\x,0,7.5);
		}

	\draw[axis] (-2,0,0) -- (8,0,0) node[left] {$x$};
	\draw[axis] (0,-1,0) -- (0,8,0) node[right] {$y$};
	\draw[axis] (0,0,-1) -- (0,0,8) node[left] {$z$};

    \draw[->,very thick,Red,dashed] (0,0,0) -- (4,0,0);
    \draw[->,very thick,Blue,dashed] (4,0,0) -- (4,3,0);
    \draw[->,very thick,Green,dashed] (4,3,0) -- (4,3,5);
    \draw[->,very thick,black] (0,0,0) -- (4,3,5)
        node[right] {\small $\vec{v} = \langle 4,3,5 \rangle$};
\end{tikzpicture}
\end{figure}

Continuing, we can see an operation that vectors can perform that
we have not discussed in the notes for linear algebra.

\begin{definition}{}{}
A \textbf{dot product} of two vectors is the sum of products of the
corresponding elements.
\end{definition}

Dot product will be quite important when we discuss attention and
transformers because it can tell us how ``related'' two vectors are.
Indeed, consider the vectors $A$ and $B$ with the angle between the
two vectors as $\theta$. The following equation holds.
\[
A \cdot B = ||A|| ||B|| \cos\theta
\]
Here, $||V||$ for $V \in \mathbb{R}^n$ represents the magnitude of the
vector, i.e. $\sqrt{\sum_{k=1}^n V_k^2}$. Let's take a look at a quick
example.
\begin{align*}
    \langle 1, 0 \rangle \cdot \langle 1, \sqrt{3} \rangle
    &= 1 \cdot 1 + 0 \cdot \sqrt{3} = 1 \\
    &= 1 \cdot 2 \cdot \cos60^\circ = 1
\end{align*}
One thing to keep in mind is that we have a special term for the
sign $||V||$.

\begin{definition}{}{}
The \textbf{vector norm} is the non-negative value that represents
the magnitude of a vector.
\end{definition}

With that in mind, let's move on to a more abstract and advanced
topics in linear algebra that are constantly used in machine learning.

\subsection{CONTINUE}






\end{document}


